Source : Insee, comptes nationaux annuels base 2014 (séries arrêtées) jusqu'en 1995 ; Eurostat et Insee base 2020 ensuite. Calculs de l'auteur (dette/dette.py).\par Lecture : de 1978 à 1995, le ratio de dette passe de 21,2 \% à 56,1 \% du PIB ; les déficits primaires n'y contribuent que pour 8,8 points, les intérêts (41,9) dépassant la croissance (-33,4). De 1995 à 2025, il passe de 57,5 \% à 115,7 \% : déficits primaires 50,5 points, intérêts 73,9, croissance -72,3, ajustements stock-flux 6,0.
